% ch-installation.tex -- requirements, build, options, test suite
\chapter{Building Yuzen}
\label{ch:installation}

Yuzen builds with CMake and has no dependency you have to install by hand. If
Eigen is not already present, CMake fetches it; if AMGCL is not present and you
asked for it, CMake fetches that too, and falls back to an Eigen-only build if
it cannot.

\section{Requirements}
\label{sec:install:requirements}

\begin{itemize}
  \item A C++17 compiler. The CI matrix covers GCC, Clang and MSVC on Linux,
        macOS and Windows.
  \item CMake $\ge$ 3.16.
  \item OpenMP, if you want the parallel build (on by default; the serial build
        is fully supported and gives the same numbers). With GCC, install
        LLVM's \code{libomp} as well (\code{apt install libomp-dev}, or the
        \code{libomp} package of a Red Hat system): the build then links it
        instead of GCC's own runtime, and runs on many cores are faster
        (\cref{sec:parallel:runtime}). \textbf{There is no MPI
        dependency and no MPI build}: Yuzen is shared-memory parallel only
        (\cref{ch:parallel}).
  \item Optionally Eigen $\ge$ 3.3 and AMGCL, if you would rather supply them
        yourself than have CMake fetch them. On Debian and Ubuntu that is
        \code{apt install libeigen3-dev}; on macOS, \code{brew install eigen};
        with vcpkg, \code{vcpkg install eigen3}.
\end{itemize}

Fetching happens only when the dependency is missing, so a machine without
network access needs Eigen installed beforehand --- and, if you want it, AMGCL.

\begin{note}[Non-ASCII file names on Windows]
  The solver handles file names as UTF-8, which is what the input file is
  and what Linux and macOS use. An MSVC build embeds a manifest
  (\file{yuzen/src/yuzen.manifest}) that makes UTF-8 the code page of the
  process, so that an input file, a \key{JobName} or a
  \key{Restart.FileName} with Japanese characters names the same file on
  Windows. That needs Windows 10 version 1903 or later; on an older Windows,
  keep those names ASCII. Paths longer than 260 characters need the Windows
  long-path setting as well.
\end{note}

\section{The build}
\label{sec:install:build}

From the top of the repository:

\begin{shell}
cmake -S . -B build -DCMAKE_BUILD_TYPE=Release
cmake --build build -j
ctest --test-dir build
\end{shell}

\code{ctest --test-dir} needs CMake 3.20 or later. An older \code{ctest}
ignores the option, finds no tests and still exits with 0, so with CMake 3.16
to 3.19 run \code{cd build \&\& ctest} instead. The whole suite takes about 25 minutes on two
cores; \code{ctest -R yuzen\_tests} runs the unit tests alone in a few seconds.

This produces two executables in \file{build/bin/}: the solver \file{yuzen} and
the test binary \file{yuzen\_tests}. Run \file{yuzen --version} and it
prints its full version --- \code{Yuzen-1.0.0}, as
\code{MAJOR.MINOR.PATCH}; with no argument it prints its usage. Input
files name only the \code{MAJOR.MINOR} series; \cref{ch:policy} says what that
distinction promises and \cref{ch:revisions} what each release changed. The
build type defaults to \val{Release}
if you do not set it; do not benchmark a \val{Debug} build, and do not draw
conclusions about convergence from one either, since the difference in rounding
is not the interesting kind.

\section{Build options}
\label{sec:install:options}

\begin{yzntable}{CMake options\label{tab:cmake-options}}{@{}L{0.30\textwidth} L{0.10\textwidth} L{0.52\textwidth}@{}}
\textsf{\bfseries Option} & \textsf{\bfseries Default} & \textsf{\bfseries Effect}\\ \midrule
\code{YUZEN\_ENABLE\_OPENMP}   & \val{ON} & OpenMP parallelization. Turning it off also removes the two extra \code{ctest} entries that check thread invariance. \\
\code{YUZEN\_WITH\_AMGCL}      & \val{ON} & Build the AMGCL backend. If the headers cannot be found or fetched, the build warns and continues without it. \\
\code{YUZEN\_STRICT\_WARNINGS} & \val{ON} & \code{-Wall -Wextra} and the implicit-fallthrough warning. \\
\code{YUZEN\_BUILD\_TESTS}     & \val{ON} & Build \file{yuzen\_tests} and register the \code{ctest} entries. \\
\code{YUZEN\_PREFER\_LIBOMP}   & \val{ON} & GCC on Linux: link LLVM's OpenMP runtime \code{libomp} instead of GCC's \code{libgomp} when it is installed; its parallel regions cost a quarter as much on 56 threads (\cref{sec:parallel:runtime}). \\
\code{YUZEN\_OPENMP\_RUNTIME}  & (empty) & The path of an OpenMP runtime to link instead of the compiler's own (\file{libomp.so}, \file{libiomp5.so}); it takes precedence over \code{YUZEN\_PREFER\_LIBOMP}. \\
\end{yzntable}

So, for example, a minimal build without AMGCL and without tests is

\begin{shell}
cmake -S . -B build -DCMAKE_BUILD_TYPE=Release \
      -DYUZEN_WITH_AMGCL=OFF -DYUZEN_BUILD_TESTS=OFF
cmake --build build -j
\end{shell}

Both settings of \code{YUZEN\_WITH\_AMGCL} are supported configurations and
both compile; the \code{\#ifdef} that distinguishes them never appears in a
header, only in \file{yuzen/src/\bk numerics/\bk LinearSolverBackend.cpp}. If you ask
for the \val{"amgcl"} backend in an input file but built without it, the run
warns and falls back to \val{"eigen"} (\cref{sec:linsolve:backends}).

\section{Installing}
\label{sec:install:install}

The solver can be run from \file{build/bin/} as it is. To put it somewhere
else,

\begin{shell}
cmake --install build --prefix $HOME/.local
\end{shell}

copies \file{yuzen} into \file{bin/} under the prefix, and \file{LICENSE},
\file{NOTICE}, \file{THIRD\_PARTY\_NOTICES.md} and \file{README.md} into
\file{share/doc/yuzen/}. Nothing else is installed: Eigen and AMGCL are
header-only and compiled in, and the scripts of \file{tools/}, the cases and
this manual are used from the source tree.

\section{The test suite}
\label{sec:install:tests}

\begin{shell}
ctest --test-dir build
\end{shell}

\code{ctest} registers seventeen entries (sixteen where SymPy is not
installed). Three run the test binary in separate
processes: once with whatever thread count the environment gives, once with
\env{OMP\_NUM\_THREADS=1} and once with \env{OMP\_NUM\_THREADS=4}. That is not
redundancy. Thread invariance is one of the claims the suite makes, and it can
only be asserted across processes (\cref{ch:reproducibility}). Twelve run the
solver itself: \code{concz\_vs\_wm} and \code{concz\_vs\_wm\_2d} check
that the two concentration models agree where they must;
\code{restart\_roundtrip} that a restart of either model, in 1-D and 2-D, from
the \file{.vts} (in \val{"CZ"} also one without the thickness and substrate) and from the \file{field*.dat}, starts from the state and the
time it was given (\cref{sec:ref-numerics:restart}), and that each \file{.vts}, read through the offsets in its header, holds what the field file of the same instant holds; \code{restart\_chain}
that a run stopped and restarted several times, from either format, gives the
same files, bit for bit, as one that ran straight through; \code{input\_guards}
that inputs which would otherwise run with a silently wrong answer --- a misspelled key, a
one-sided \val{"Periodic"}, a \val{"Periodic"} edge on a rotating film, a
pinned edge facing a film of full thickness ---
are refused with the reason; \code{error\_paths} that what a new user, a full
disk or a job scheduler does to a run --- a wrong command line, input file or
result directory, a missing or truncated restart file, a file size limit or a
full device, \code{SIGINT}, \code{SIGTERM} or \code{SIGKILL} in the middle
of a run, paths with spaces and Japanese characters --- stops it with the path
and the reason and leaves no file cut short, and that a run interrupted and
continued is, bit for bit, the run made straight through
(\cref{sec:output:failures}); \code{equivalence} that the same problem
posed in two ways gives the same answer (\cref{sec:verification:equivalence});
\code{analytic} that the problems with a closed-form answer get it, at
the order the scheme promises (\cref{sec:verification:analytic}); \code{mms}
that a manufactured solution with a large amplitude, with every term of each
model on, converges at second order in space and time
(\cref{sec:verification:mms}); \code{advection} that the convection schemes
keep a concentration step within its bounds where they promise to, and
converge on a smooth one at their order, and that the well-mixed
concentration converges in 2-D at second order with convection and
diffusion together (\cref{sec:verification:schemes});
\code{features} that the probes, the
steady-state, dry-out and \key{MaxSteps} stops and the growth of an adaptive
time step do what they say (\cref{sec:verification:features}); and
\code{proptest} that random valid inputs end cleanly, keep the film positive
and the volume and the solute, and are invariant where they must be
(\cref{sec:verification:fuzzing}). One runs no
solver: \code{doc\_keys} reads every key the parser reads, with its type,
default, range and accepted values, holds the key tables of
\cref{ch:ref-domain,ch:ref-physics,ch:ref-bcout,ch:ref-numerics}, the TOML
listings of this manual and every file in \file{cases/} to it, and has the
binary read each case with \code{--check-input}
(\cref{sec:inputformat:unknown}). The last,
\code{mms\_source\_current}, is registered only where SymPy is installed: it
checks that the generated source of the manufactured solution is up to date.
All but the three test-binary entries need a Python interpreter and are left
out without one.

Running the test binary directly prints the counts:

\begin{plain}
$ ./build/bin/yuzen_tests
[doctest] test cases:    128 |    128 passed | 0 failed | 0 skipped
[doctest] assertions: 183297 | 183297 passed | 0 failed |
\end{plain}

These are not smoke tests. They assert against analytic solutions, conservation
laws and convergence orders; \cref{sec:verification:suite} lists what each of
the seventeen test files fixes.

\section{Repository layout}
\label{sec:install:layout}

\begin{plain}
yuzen/src/grid/        structured grid and fields (staggered, type-tagged)
yuzen/src/numerics/    vectors, CRS matrices, linear solvers, time integrators
yuzen/src/model/       the physics: pressure terms, equations, model drivers
yuzen/src/IO/          input parsing, output, restart, probes
yuzen/src/util/        logging, paths, job information
yuzen/src/commonFunctions/   shape functions and shape generation
yuzen/tests/           the test suite
cases/                 example and regression input files
tools/                 helper scripts
docs/                  theory, numerics, validation, and this manual
third_party/           doctest, toml++
\end{plain}

\Cref{ch:layout} goes through this in more detail, for readers who intend to
modify the solver.

\section{Continuous integration}
\label{sec:install:ci}

\file{.github/workflows/ci.yml} builds and tests on Ubuntu with GCC and Clang
(and with GCC 11, the system compiler of RHEL 9 and of many clusters),
on Windows with MSVC, and on macOS with Clang, all \emph{without} AMGCL so
that the plain configuration is the one kept honest, and once more on Ubuntu
with GCC and AMGCL, so that the code path taken only with AMGCL is built and
tested too.

\section{If the build fails}
\label{sec:install:trouble}

\begin{description}
  \item[Eigen is fetched every time.] CMake only fetches when
    \code{find\_package(Eigen3 3.3 QUIET NO\_MODULE)} fails. Install Eigen
    system-wide, or point \code{Eigen3\_DIR} at it, and the fetch stops.
  \item[AMGCL cannot be fetched.] The build warns, sets
    \code{YUZEN\_WITH\_AMGCL} to \val{OFF} and continues. Nothing else is lost;
    you simply have one backend, \val{"eigen"}, instead of two.
  \item[OpenMP not found on macOS.] Apple Clang needs \code{libomp}
    (\code{brew install libomp}). If you would rather not, configure with
    \code{-DYUZEN\_ENABLE\_OPENMP=OFF}: the serial build produces exactly the
    same numbers.
  \item[Warnings you do not want.] Configure with
    \code{-DYUZEN\_STRICT\_WARNINGS=OFF}. Please do not make this the default
    for development work; the warning set was turned on deliberately.
\end{description}
